\documentclass[11pt,a4paper]{article}
\usepackage[margin=25mm]{geometry}
\usepackage{amsmath,amssymb,amsthm}
\usepackage{booktabs}
\usepackage{hyperref}
\usepackage{enumitem}
\setlist{nosep}

\newtheorem{proposition}{Proposition}
\newtheorem{definition}{Definition}
\theoremstyle{remark}
\newtheorem{remark}{Remark}
\newcommand{\N}{\mathcal N}
\newcommand{\E}{\mathbb E}
\newcommand{\R}{\mathbb R}
\newcommand{\Cov}{\operatorname{Cov}}
\newcommand{\pos}[1]{(#1)_+}

\title{The mixture form of the Version-4 basket model:\\ model, likelihood, sampling and training}
\author{Ajit}
\date{27 October 2026}

\begin{document}
\maketitle

\begin{abstract}
This document states the model completely and derives everything that is computed from it. Section~1
fixes notation. Section~2 defines the model: a prior over a latent taste $z$, an exact law of the basket
given $z$, the marginal law of the basket, and the posterior of $z$ given a basket. Section~3 shows how
the law given $z$ is normalised exactly, and Section~4 how a basket is drawn from it exactly, as a chain of
three conditional draws. Section~5 explains how the likelihood of an observed basket is estimated, and
Section~6 how the parameters are trained. Section~7 covers generation and counterfactuals, Section~8
compares with the earlier form of the model, and Section~9 gives the evidence so far.
\end{abstract}

\tableofcontents

%======================================================================
\section{Notation}

\begin{itemize}
\item A \emph{trip} $t$ is one store visit by household $h=h(t)$ at store $s(t)$ in week $w(t)$. Its context
  $x_t$ holds prices, discounts, display and mailer flags, the week, the store and the household's purchase
  history before the trip.
\item $\mathcal J_t$ is the set of products offered on trip $t$. Every product $j$ belongs to one category
  $c(j)\in\{1,\dots,C\}$; $\mathcal J_{tc}=\{j\in\mathcal J_t: c(j)=c\}$. On the Dunnhumby commodity data
  $C=187$ and there are 5{,}455 products.
\item A \emph{basket} is a non-empty set $S\subseteq\mathcal J_t$. We write
  \[
  n_c(S)=|S\cap\mathcal J_{tc}|,\qquad n(S)=\sum_c n_c(S),\qquad m(S)=\#\{c: n_c(S)>0\}
  \]
  for the number of products in category $c$, the basket size and the number of categories present.
\item Support: $\mathcal S_t=\{S\subseteq\mathcal J_t:\ 1\le n(S)\le n_{\max},\ n_c(S)\le R\ \forall c\}$
  with $n_{\max}=120$ and $R=40$. All laws below are on $\mathcal S_t$.
\item $r$ is the taste dimension ($r=16$). $\N(z;\mu,\Sigma)$ is the Gaussian density on $\R^r$.
  $\pos{a}=\max(a,0)$, $\binom{k}{2}=k(k-1)/2$.
\end{itemize}
The trip index $t$ is dropped when only one trip is involved.

%======================================================================
\section{The model}

\subsection{Parameters}

\begin{center}\small
\begin{tabular}{lp{2cm}p{7.2cm}l}\toprule
symbol & per & meaning & trained?\\\midrule
$\lambda_j$ & product & popularity intercept & yes\\
$\phi_j\in\R^r$ & product & taste loading of product $j$ & yes\\
$\psi_c\in\R^r$ & category & taste loading of an \emph{extra} product of category $c$ & yes\\
$\rho_c$ & category & penalty per pair of products within $c$ & yes\\
$f_c(2),f_c(3),f_c(\ge4)$ & category & free penalty on holding 2, 3, $\ge4$ products of $c$ & yes\\
$\rho_0(1..n_{\max})$ & count & size potential on $m$, the number of categories present & yes\\
rest of $b_{tj}$ & & context utility, \eqref{eq:b} & yes${}^{*}$\\\bottomrule
\end{tabular}
\end{center}
The fit (\texttt{scripts/fit\_mixture.py}) starts from the exact additive base model ($\phi=0$) and trains
every parameter except the within-category competition weights in $b_{tj}$, which the price instrument
re-estimates once per epoch. ${}^{*}$The real-data fit of Section~\ref{sec:evidence} is older: it started from
the earlier form's checkpoint \texttt{runs/commodity/\allowbreak v4\_comp\_promo.pt} and kept the rest of
$b_{tj}$ at its values.

\subsection{Context utility}

For product $j$ on trip $t$,
\begin{align}
b_{tj}=\lambda_j &+\theta_h\!\cdot\!\alpha_j+\theta_h\!\cdot\!u_{c(j)}
 -g_{hj}\big[\bar p_t+\kappa(p_{tj}-\bar p_t)\big]-g_{hj}(\kappa_d-\kappa)\,d_{tj}-\gamma_d\,d_{tj}\nonumber\\
&+w^{\rm dsp}_j D_{tj}+w^{\rm mlr}_j M_{tj}+\mu_j\!\cdot\!\delta_{w(t)}+\zeta_j\!\cdot\!\xi_{s(t)}
 +\eta_t(j)-o_t+\iota_{t,c(j)} .
\label{eq:b}
\end{align}
$\theta_h\cdot\alpha_j$ and $\theta_h\cdot u_c$ are household tastes (product and category);
$p_{tj}$ is the log price deviation, $\bar p_t$ its mean over the trip's assortment, $d_{tj}$ the discount
depth, $g_{hj}\ge0$ the household-by-product price sensitivity, $\kappa,\kappa_d>0$ the relative-price and
discount multipliers, $\gamma_d$ a common discount slope; $D,M$ are display and mailer flags;
$\mu_j\cdot\delta_w$ and $\zeta_j\cdot\xi_s$ are season and store terms; $\eta_t(j)$ is the purchase-history
term; $o_t$ is a trip-level offset (the mean of $\eta_t$ over the assortment plus a history-size term), the
same for every product of the trip; and
\[
\iota_{tc}=(\ell_c-1)\big[\mathrm{IV}_{tc}-\mathrm{IV}^{\rm reg}_{tc}\big],\qquad
\mathrm{IV}_{tc}=\log\!\!\sum_{j\in\mathcal J_{tc}}\!\! e^{b^{\circ}_{tj}},\quad \ell_c\in(0,1),
\]
is the within-category competition shift ($b^\circ$ is \eqref{eq:b} without $\iota$ and $o_t$, and ``reg'' the same with
price deviation, discount, display and mailer set to zero). With $\ell_c$ near 0 a price or promotion change
in $c$ moves purchases between products of $c$ rather than raising the category's purchase rate. In
this document $b_{tj}$ is simply a known number per (trip, product): it does not depend on $z$ or on the
basket.

\subsection{Prior}

Each trip has a latent taste
\begin{equation}
z_t\sim\N(0,I_r),\qquad\text{independently across trips.}
\label{eq:prior}
\end{equation}
$z_t$ describes what the shopper is in the mood for on this trip (for example ``baking'' or ``baby
products''). It is not tied to the household: two trips of the same household draw independent tastes;
persistent household preferences are in $\theta_h$ inside $b_{tj}$.

\subsection{The law of the basket given the taste}

\begin{definition}[conditional law]
Given $z$, the probability of basket $S\in\mathcal S_t$ is
\begin{equation}
p(S\mid z)=\frac{\exp E_z(S)}{Z_z},\qquad Z_z=\sum_{S'\in\mathcal S_t}\exp E_z(S'),
\label{eq:cond}
\end{equation}
with
\begin{align}
E_z(S)=\sum_{j\in S}u_j(z)
-\sum_{c}h_c\big(n_c(S),z\big)-\rho_0\big(m(S)\big),
\label{eq:Ez}
\end{align}
\begin{align}
u_j(z)&=b_j-\tfrac12\lVert\phi_j\rVert^2+z\cdot\phi_j, \label{eq:u}\\
h_c(k,z)&=\rho_c\tbinom{k}{2}+f_c(k)+\pos{k-1}\big(z\cdot\psi_c-\tfrac12\lVert\psi_c\rVert^2\big),
\qquad h_c(0,z)=0. \label{eq:h}
\end{align}
Here $f_c(0)=f_c(1)=0$, $f_c(k)=f_c(\ge4)$ for $k\ge4$, and $\rho_0(1)=0$ fixes the free constant.
\end{definition}

Reading \eqref{eq:Ez}: each product contributes its utility $u_j(z)$; a taste $z$ raises the products whose
loading $\phi_j$ points the same way. Within a category, every additional product beyond the first pays
$\rho_c$ per pair and $f_c$, and its taste contribution is corrected by $\psi_c$ (a category that is liked
need not be bought many times). The size term charges the number of categories present.

\paragraph{The basket field.} Collect the $z$-dependence of $E_z(S)$:
\begin{equation}
E_z(S)=a(S)+z\cdot s(S),\qquad
s(S)=\sum_{j\in S}\phi_j-\sum_c\pos{n_c(S)-1}\,\psi_c\in\R^r,
\label{eq:affine}
\end{equation}
\[
a(S)=\sum_{j\in S}\big(b_j-\tfrac12\lVert\phi_j\rVert^2\big)-\sum_c\Big[\rho_c\tbinom{n_c}{2}+f_c(n_c)
-\tfrac12\pos{n_c-1}\lVert\psi_c\rVert^2\Big]-\rho_0(m).
\]
$a(S)$ does not involve $z$. Equation~\eqref{eq:affine}, \emph{the energy is affine in $z$}, is the one fact
every property below uses. $s(S)$ is called the basket field: the taste direction the basket points in.

\paragraph{Why $-\tfrac12\lVert\phi_j\rVert^2$?} It makes the expected boost of product $j$ under the prior
zero on the multiplicative scale: $\E_{z\sim\N(0,I)}e^{z\cdot\phi_j-\frac12\lVert\phi_j\rVert^2}=1$. It is a
choice of parameterisation (it could be absorbed into $\lambda_j$), kept so that the earlier form and the
mixture form share their parameters.

\subsection{The marginal law: the model}

\begin{definition}[mixture form]
The probability that trip $t$ produces basket $S$ is
\begin{equation}
p(S)=\int_{\R^r}\N(z;0,I)\,p(S\mid z)\,dz .
\label{eq:marginal}
\end{equation}
\end{definition}
This is a continuous mixture: one basket law $p(\cdot\mid z)$ for every taste $z$, weighted by the prior
density of that taste. Equivalently, the data are generated hierarchically:
\[
z_t\sim\N(0,I),\qquad S_t\mid z_t\sim p(\cdot\mid z_t).
\]
Products bought together are explained by a shared taste. The number of components is not a choice:
there is one per $z$, they all share the formula \eqref{eq:cond}, and the parameter count does not depend
on it. The modelling choice is the dimension $r$ and the Gaussian shape of the prior. $p(S)$ is a proper
probability: $\sum_S p(S)=\int\N(z)\sum_S p(S\mid z)\,dz=1$.

\subsection{The joint law and the posterior}

The joint density of taste and basket is $p(z,S)=\N(z;0,I)\,p(S\mid z)$. By Bayes' rule the posterior of
the taste given an observed basket is
\begin{equation}
p(z\mid S)=\frac{\N(z;0,I)\,p(S\mid z)}{p(S)},\qquad
\log p(z\mid S)=-\tfrac12\lVert z\rVert^2+z\cdot s(S)-\log Z_z+\text{const}(S),
\label{eq:post}
\end{equation}
where the second form uses \eqref{eq:affine}; the constant collects $a(S)$, $\log p(S)$ and the Gaussian
constant. The posterior is the main computational object: the likelihood $p(S)$ is its normaliser, and
training averages over it.

%======================================================================
\section{Computing the conditional law exactly}
\label{sec:dp}

$Z_z$ is a sum over all baskets, astronomically many. It is computed exactly because, given $z$, the
weight $e^{E_z(S)}$ factorises over categories except for the size term, which depends on one integer.

\paragraph{Elementary symmetric polynomials.} Write $w_j=e^{u_j(z)}$. For category $c$ and $k\ge0$,
\[
e_{c,k}=\sum_{A\subseteq\mathcal J_c,\ |A|=k}\ \prod_{j\in A}w_j ,
\]
the total product weight of all ways to choose $k$ products from $c$. These follow from the recursion
over the products of $c$ in a fixed order: $e^{(p)}_k=e^{(p-1)}_k+w_p\,e^{(p-1)}_{k-1}$, $e^{(0)}_0=1$.

\paragraph{Category polynomials.} The total weight of category $c$ holding exactly $k$ products is
\begin{equation}
G_c(k)=e_{c,k}\,e^{-h_c(k,z)},\qquad G_c(0)=1 .
\label{eq:G}
\end{equation}

\paragraph{The size table.} Let $F_c(n,m)$ be the total weight of all choices in categories $1,\dots,c$
that hold $n$ products in $m$ present categories:
\begin{equation}
F_0(n,m)=[n=0,m=0],\qquad
F_c(n,m)=F_{c-1}(n,m)+\sum_{k=1}^{\min(R,n)}F_{c-1}(n-k,m-1)\,G_c(k).
\label{eq:F}
\end{equation}
Equivalently $F_C(n,m)$ is the coefficient of $t^n y^m$ in $\prod_c\big(1+y\sum_{k\ge1}G_c(k)t^k\big)$.

\begin{proposition}[exact normaliser]
\begin{equation}
Z_z=\sum_{n=1}^{n_{\max}}\sum_{m=1}^{n}F_C(n,m)\,e^{-\rho_0(m)} .
\label{eq:Z}
\end{equation}
\end{proposition}
\begin{proof}
A basket is a choice, for each category, of a subset of size $n_c\le R$. Its weight $e^{E_z(S)}$ is
$e^{-\rho_0(m)}\prod_c\big(\prod_{j\in S_c}w_j\big)e^{-h_c(n_c,z)}$. Summing over the subsets with the counts
$(n_c)$ fixed gives $e^{-\rho_0(m)}\prod_c G_c(n_c)$; \eqref{eq:F} sums $\prod_cG_c(n_c)$ over all count
vectors with total $n$ and $m$ non-zero entries.
\end{proof}

The cost is $O(\sum_c|\mathcal J_c|R)$ for the polynomials and $O(C\,n_{\max}^2R)$ for the table, per trip
and per $z$. Numerically every row $n$ of the table carries its own power-of-two exponent, because weights
of different sizes differ by far more than a double can hold. The gradients follow by differentiating
\eqref{eq:Z}:
\begin{equation}
\nabla_z\log Z_z=\E[s(S)\mid z],\qquad \nabla^2_z\log Z_z=\Cov[s(S)\mid z],\qquad
\frac{\partial\log Z_z}{\partial u_j}=P(j\in S\mid z),
\label{eq:grads}
\end{equation}
the first two because $\log Z_z=\log\sum_Se^{a(S)+z\cdot s(S)}$. All three are exact, computed by
reverse-mode differentiation through \eqref{eq:G}--\eqref{eq:Z}. Likewise, for any model parameter
$\vartheta$, $\nabla_\vartheta\log Z_z=\E_{S\sim p(\cdot\mid z)}[\nabla_\vartheta E_z(S)]$.

%======================================================================
\section{Sampling}
\label{sec:sampling}

\subsection{A basket given a taste: three conditional draws}

Given $z$, a basket is drawn exactly from \eqref{eq:cond} by the following chain of conditional draws,
using the tables of Section~\ref{sec:dp}.

\begin{enumerate}[label=\textbf{Step \Alph*.},leftmargin=*]
\item \textbf{Size and number of categories.} Draw $(n,m)$ with
  \begin{equation}
  P(n,m\mid z)=\frac{F_C(n,m)\,e^{-\rho_0(m)}}{Z_z}.
  \label{eq:stepA}
  \end{equation}
\item \textbf{Count per category}, going backwards $c=C,C-1,\dots,1$. Let $(n',m')$ be the products and
  categories still to place in categories $1..c$ (initially $(n,m)$). Draw $n_c=k$ with
  \begin{equation}
  P(n_c=k\mid n',m',z)=\frac{F_{c-1}\big(n'-k,\ m'-[k>0]\big)\,G_c(k)}{F_c(n',m')},\qquad 0\le k\le\min(R,n'),
  \label{eq:stepB}
  \end{equation}
  then set $n'\leftarrow n'-k$, $m'\leftarrow m'-[k>0]$. The probabilities sum to one by \eqref{eq:F}, and at the
  end $n'=m'=0$.
\item \textbf{Products within each category.} For every category with $n_c=k>0$, draw a subset $A$ of
  $\mathcal J_c$ of size $k$ with
  \begin{equation}
  P(A\mid n_c=k,z)=\frac{\prod_{j\in A}w_j}{e_{c,k}}.
  \label{eq:stepC}
  \end{equation}
  This is done sequentially: go through the products $p=L,L-1,\dots,1$ of the category with $q$ still to
  choose, and include $p$ with probability $w_p\,e^{(p-1)}_{q-1}/e^{(p)}_q$ (after which $q\leftarrow q-1$).
\end{enumerate}

\begin{proposition}[exactness]
The basket produced by Steps A--C has law $p(S\mid z)$.
\end{proposition}
\begin{proof}
Multiply the probabilities. In Step B the denominators telescope against the numerators of the next
category, and Step A's numerator cancels the first denominator, so
\[
P(n,m)\prod_{c}P(n_c\mid\cdot)=\frac{e^{-\rho_0(m)}}{Z_z}\,F_0(0,0)\prod_cG_c(n_c)
=\frac{e^{-\rho_0(m)}}{Z_z}\prod_c e_{c,n_c}e^{-h_c(n_c,z)} .
\]
In Step C the inclusion probabilities telescope in the same way, because $e^{(p)}_q=e^{(p-1)}_q+w_pe^{(p-1)}_{q-1}$,
and their product is \eqref{eq:stepC}. Multiplying by $\prod_c\prod_{j\in A_c}w_j/e_{c,n_c}$ gives
\[
\frac{e^{-\rho_0(m)}}{Z_z}\prod_c e^{-h_c(n_c,z)}\prod_{j\in S}w_j=\frac{e^{E_z(S)}}{Z_z}.
\]
\end{proof}

Sampling is therefore hierarchical: size and categories first, then the count in each category given what
is left, then the products given the count. Each conditional is the exact conditional of $p(S\mid z)$.

\subsection{A basket from the model}

\begin{proposition}[exact draws from the model]
Draw $z\sim\N(0,I)$, then $S$ by Steps A--C given $z$. Then $S\sim p(S)$ of \eqref{eq:marginal}.
\end{proposition}
\begin{proof}
The pair $(z,S)$ has density $\N(z;0,I)\,p(S\mid z)$; its $S$-marginal is \eqref{eq:marginal}.
\end{proof}
No Markov chain is involved, so there is no burn-in, no mixing question and no correlation between draws.
The full hierarchy of a generated basket is therefore
\[
z\ \to\ (n,m)\ \to\ (n_1,\dots,n_C)\ \to\ \text{products}.
\]

\paragraph{Truncation for speed.} The sampler first builds the table only up to a per-trip cap
$K<n_{\max}$ chosen so that baskets larger than $K$ carry at most $10^{-8}$ of $Z_z$ (a proven bound); if
the bound is not met the table is rebuilt to $n_{\max}$. The draw is exact up to that $10^{-8}$.

%======================================================================
\section{The likelihood of an observed basket}
\label{sec:lik}

$p(S)$ in \eqref{eq:marginal} is an $r=16$-dimensional integral with no closed form. It is estimated by
importance sampling over $z$, which works well because the integrand has one peak.

\subsection{The posterior has one peak}

\begin{proposition}[strong log-concavity]\label{prop:concave}
For every basket $S$, $\log p(z\mid S)$ is 1-strongly concave in $z$:
\[
\nabla_z\log p(z\mid S)=s(S)-z-\E[s\mid z],\qquad \nabla^2_z\log p(z\mid S)=-\big(I+\Cov[s\mid z]\big)\preceq-I .
\]
Hence it has exactly one maximiser $\hat z(S)$, the solution of $\hat z=s(S)-\E[s\mid\hat z]$, and
$\log p(z\mid S)\le\log p(\hat z\mid S)-\tfrac12\lVert z-\hat z\rVert^2$ for all $z$.
\end{proposition}
\begin{proof}
Differentiate \eqref{eq:post} using \eqref{eq:grads}. A covariance is positive semi-definite, so the
Hessian is at most $-I$. A strongly concave function has a unique maximiser, and the quadratic bound is
Taylor's theorem with the Hessian bound.
\end{proof}

\paragraph{Finding the peak.} Gradient ascent $z\leftarrow z+\eta\,\nabla_z\log p(z\mid S)$, with the step
$\eta\le1$ halved until $\log p(z\mid S)$ rises by at least $\tfrac14\eta\lVert\nabla\rVert^2$ (Armijo), so every
iteration increases it. Each iteration needs $\E[s\mid z]$, one exact gradient of $\log Z_z$. Validation
runs 30 iterations from $z=0$; training runs 8 starting from the trip's peak at its previous visit.

\subsection{The proposal}

Around the peak the posterior is close to the Gaussian with covariance equal to the inverse negative Hessian
(Laplace approximation),
\[
\Sigma(S)=\big(I+\Cov[s\mid\hat z]\big)^{-1},
\]
where $\Cov[s\mid\hat z]$ is the sample covariance of $s(S')$ over 32 exact draws $S'\sim p(\cdot\mid\hat z)$
(Section~\ref{sec:sampling}). The proposal is
\begin{equation}
q(z\mid S)=0.8\,\N\big(z;\hat z,\ 1.2^2\Sigma(S)\big)+0.2\,\N\big(z;\hat z,\ I\big).
\label{eq:q}
\end{equation}
The first part follows the posterior's shape, slightly widened; the second has unit covariance, which by
Proposition~\ref{prop:concave} is at least as wide as the posterior in every direction.

\subsection{The estimator}

Draw $z_1,\dots,z_K\sim q(\cdot\mid S)$ independently and set
\begin{equation}
\omega_k=\frac{\N(z_k;0,I)\,p(S\mid z_k)}{q(z_k\mid S)},\qquad
\hat p_K(S)=\frac1K\sum_{k=1}^K\omega_k .
\label{eq:is}
\end{equation}
Each $p(S\mid z_k)$ is exact by \eqref{eq:cond} and \eqref{eq:Z}.

\begin{proposition}[properties of the estimate]\label{prop:is}
\begin{enumerate}
\item $\E\,\hat p_K(S)=p(S)$ (unbiased).
\item $\E\log\hat p_K(S)\le\E\log\hat p_{K+1}(S)\le\log p(S)$, and $\E\log\hat p_K(S)\to\log p(S)$ as $K\to\infty$.
\item Every weight is bounded: $\omega_k\le(2\pi)^{r/2}\,p(S)\,p(\hat z\mid S)/0.2$.
\end{enumerate}
\end{proposition}
\begin{proof}
(1) $\E_q[\N p(S\mid z)/q]=\int\N(z)p(S\mid z)\,dz$. (2) The upper bound is Jensen's inequality; monotonicity
in $K$ and convergence are the importance-weighted bound of Burda, Grosse and Salakhutdinov (2016).
(3) $\N(z;0,I)p(S\mid z)=p(S)p(z\mid S)\le p(S)p(\hat z\mid S)e^{-\frac12\lVert z-\hat z\rVert^2}$ by Proposition~\ref{prop:concave}, and
$q(z)\ge0.2\,(2\pi)^{-r/2}e^{-\frac12\lVert z-\hat z\rVert^2}$.
\end{proof}
So the reported log-likelihood can understate the model but on average never overstates it, and no single
draw can dominate without bound. The effective sample size
\[
\mathrm{ESS}=\frac{(\sum_k\omega_k)^2}{K\sum_k\omega_k^2}\in(0,1]
\]
is reported as a quality check. Validation uses $K=256$.

\paragraph{Posterior averages.} The normalised weights $\bar\omega_k=\omega_k/\sum_l\omega_l$ estimate
averages over the posterior, $\E[g(z)\mid S]\approx\sum_k\bar\omega_kg(z_k)$. This estimate is consistent as
$K\to\infty$ but biased for small $K$: it leans towards the proposal. Section~\ref{sec:whiten} shows where
this matters.

%======================================================================
\section{Training}
\label{sec:train}

\subsection{Objective}

Let $\vartheta$ be the trained parameters (Section~2.1) and $t=1,\dots,N$ the training trips
($N=158{,}983$). The target is the penalised log-likelihood
\[
\frac1N\sum_{t=1}^N\log p_\vartheta(S_t)\ -\ \frac{1}{N}\Pi(\vartheta),\qquad
\Pi=\tfrac12\Big[20\sum_c\lVert\psi_c\rVert^2+200\sum_j\lVert\phi_j-\psi_{c(j)}\rVert^2+0.2\sum_c\rho_c^2+2\sum_{c,i}f_{c,i}^2\Big],
\]
a Gaussian prior that keeps each product's loading near its category's, plus the Gaussian priors on the
rest of $b_{tj}$: price parameters centred at the base model's values, the household taste per category
and the history terms at zero, and a pooling penalty on the three product-by-context factorisations
(\texttt{v4/training.py}). Since $\log p_\vartheta(S_t)$ is not
available, training maximises the $K$-sample bound of Proposition~\ref{prop:is},
\begin{equation}
\mathcal L_K(\vartheta)=\frac1N\sum_t\E_{z_{t,1..K}\sim q_t}\Big[\log\frac1K\sum_{k=1}^K
\frac{\N(z_{tk};0,I)\,p_\vartheta(S_t\mid z_{tk})}{q_t(z_{tk})}\Big]-\frac{\Pi(\vartheta)}N ,
\label{eq:obj}
\end{equation}
with $K=8$. The proposal $q_t$ is built from the current parameters but treated as fixed when
differentiating: the bound holds for any proposal, so this is legitimate.

\subsection{Gradient}

For fixed draws, $\nabla_\vartheta\log\hat p_K(S)=\sum_k\bar\omega_k\nabla_\vartheta\log p_\vartheta(S\mid z_k)$, and
\begin{equation}
\nabla_\vartheta\log p_\vartheta(S\mid z)=\nabla_\vartheta E_z(S)-\E_{S'\sim p(\cdot\mid z)}\big[\nabla_\vartheta E_z(S')\big].
\label{eq:grad}
\end{equation}
The first term is the observed basket's; the second is the model's expectation given the same taste, exact
by reverse-mode differentiation through the dynamic program of Section~\ref{sec:dp}. No Markov chain and no
model samples enter the gradient. The gradient is the posterior-weighted average of ``observed minus
expected given the taste''.

\subsection{One training step}

For a batch of 256 trips:
\begin{enumerate}
\item Find each trip's posterior peak $\hat z_t$ (8 Armijo iterations, warm-started).
\item Draw 32 baskets $S'\sim p(\cdot\mid\hat z_t)$ exactly and form $\Sigma(S_t)$ and the proposal \eqref{eq:q}.
\item Draw $z_{t1},\dots,z_{t8}\sim q_t$.
\item Compute $\log\hat p_8(S_t)$ with gradients, add the penalty, and take one Adam step (gradient norm
  clipped at 10). Learning rates: $3\times10^{-3}$ for $\phi,\psi$; $0.05$, $0.03$, $0.006$, $0.008$ for
  $\rho_0$, $\rho_c$, $f$, $\lambda$.
\item Apply the bounds: $\rho_c\ge-1.5$ and the floors on $f_c$ and the category reward.
\end{enumerate}
Every 30 steps two extra updates follow (Sections~\ref{sec:mm} and \ref{sec:whiten}). An epoch is 622 steps;
the fit runs 2 epochs.

\subsection{Moment matching}
\label{sec:mm}

At a maximum of the likelihood the gradient in $\lambda_j$ vanishes, and by \eqref{eq:grad} that gradient
is ``observed rate of $j$ minus model rate of $j$'' (posterior-weighted). The same holds for $\rho_0(m)$
(frequency of $m$ categories present) and $\rho_c$ (mean $\binom{n_c}{2}$). Adam approaches this slowly for
rarely bought products, so every 30 steps each of these parameters takes the step
\[
\lambda_j\leftarrow\lambda_j+0.35\,\mathrm{clip}\Big(\frac{\text{data rate}_j-\text{model rate}_j}{\text{data rate}_j},\pm1\Big),
\]
with the rates averaged over the last 30 steps (and factors $0.5$ for $\rho_0,\rho_c$). For small errors
$(\text{data}-\text{model})/\text{data}\approx\log(\text{data}/\text{model})$, which is the change in $\lambda_j$
that makes a Poisson rate match. This changes the speed of convergence, not its target.

\subsection{Whitening (parameter-expanded EM)}
\label{sec:whiten}

\paragraph{The idea.} Suppose the prior were $\N(0,\Sigma_0)$ with $\Sigma_0$ a free parameter. The EM update
of $\Sigma_0$ is the average posterior second moment
\begin{equation}
\Sigma_0^{\rm new}=\frac1N\sum_t\E\big[z z^\top\mid S_t\big].
\label{eq:em}
\end{equation}
A model with prior $\N(0,\Sigma_0)$ is the same law as one with prior $\N(0,I)$ and loadings rescaled: put
$z=A\tilde z$ with $A=\Sigma_0^{1/2}$ (symmetric), $\tilde z\sim\N(0,I)$. Then $z\cdot\phi_j=\tilde z\cdot A\phi_j$ and
$z\cdot\psi_c=\tilde z\cdot A\psi_c$. So each whitening step sets
\[
\phi_j\leftarrow A\phi_j,\qquad\psi_c\leftarrow A\psi_c,\qquad
\lambda_j\leftarrow\lambda_j+\tfrac12\big(\lVert A\phi_j\rVert^2-\lVert\phi_j\rVert^2\big),
\]
the last so that $u_j$ in \eqref{eq:u} is unchanged. This is an EM step for the scale and orientation of the
tastes, which gradient steps on $\phi$ alone find slowly. If the model is right, $\Sigma_0^{\rm new}=I$ and the step
does nothing.

\paragraph{Not exactly the same law.} The term $\tfrac12\pos{n_c-1}\lVert\psi_c\rVert^2$ in \eqref{eq:h} becomes
$\tfrac12\pos{n_c-1}\lVert A\psi_c\rVert^2$ and is not compensated, so the step is an approximate EM step.

\paragraph{Estimating \eqref{eq:em}.} The posterior second moments are estimated with normalised importance
weights (Section~5.3) on 256 fixed training trips with $K=256$. An earlier version reused the training
draws ($K=8$); that estimate is biased upward, because with few draws the normalised weights lean toward the
proposal, which is wider than the posterior. Measured on the same 1{,}024 trips and the same model:
\begin{center}\small
\begin{tabular}{lcc}\toprule
estimate of $\frac1{rN}\operatorname{tr}\sum_t\E[zz^\top\mid S_t]$ & $K=8$ & $K=256$\\\midrule
value (1 = model correctly scaled) & 1.056 & 0.990\\\bottomrule
\end{tabular}
\end{center}
With $K=8$ every whitening enlarged the tastes by about 3\% that the data did not ask for; over many rounds
generated baskets grew from 7.66 to 8.23 products (data 7.26). With $K=256$ this stops.

\subsection{Validation}

Every 120 steps: $\log\hat p_{256}$ on 64 held-out trips (the fit's validation split; test weeks are not
used), and the size distribution of exact draws (Section~4.2) for 1{,}024 held-out trips against the data.
The checkpoint with the best validation log-likelihood is kept.

%======================================================================
\section{Generation and counterfactuals}

\paragraph{Generation.} For trip $t$: build $b_{tj}$ from the context, draw $z\sim\N(0,I)$, then Steps A--C.
Repeating gives the distribution of baskets for that trip. Incidence has an exact inner part:
\[
P(j\in S)=\E_{z\sim\N(0,I)}\big[P(j\in S\mid z)\big],\qquad P(j\in S\mid z)=\partial\log Z_z/\partial u_j\ \text{from \eqref{eq:grads}},
\]
so only the outer average over $z$ is Monte Carlo.

\paragraph{Counterfactuals.} A price or promotion change alters $b_{tj}$ (and $\iota_{tc}$) only. Draw $z$ once,
then draw $S$ under the old and the new context with the same $z$: the same shopper's taste before and after
the change. In this form the distribution of tastes is fixed by the prior and does not respond to prices;
prices change what each taste buys. Effects on whether a trip happens at all are outside the model.

\paragraph{Conditioning on part of a basket.} For recommendations given a partly observed basket $T$,
average over $p(z\mid T\subseteq S)\propto\N(z)P(T\subseteq S\mid z)$. Unlike Proposition~\ref{prop:concave} this posterior is
not guaranteed to have one peak, because $\log P(T\subseteq S\mid z)$ is a difference of two convex functions.

%======================================================================
\section{Relation to the earlier form}

The earlier Version-4 law was $p_{\rm old}(S)=e^{a(S)+\frac12\lVert s(S)\rVert^2}/Z$ with the same $a$ and $s$. Using
$e^{\frac12\lVert s\rVert^2}=\int\N(z;0,I)e^{z\cdot s}dz$,
\[
p_{\rm old}(S)=\int\underbrace{\frac{\N(z;0,I)\,Z_z}{Z}}_{\text{weight of taste }z}\ p(S\mid z)\,dz .
\]
Both forms use the same $p(S\mid z)$; they differ in the weight of a taste. In the earlier form a taste is
weighted by $\N(z)Z_z$, and $Z_z$ grows enormously for tastes that make large baskets attractive. Its logarithm
$-\frac12\lVert z\rVert^2+\log Z_z$ is concave plus convex and can have several peaks. On the Dunnhumby refit
one peak, far from the data, made 120-product baskets carry most of the probability. Training (Gibbs chains
started at the data) and validation (annealed chains) both missed it: the reported held-out log-likelihood was
$-38.8$ per trip, and $-53.8$ once that peak was included. In the mixture form the weight is $\N(z)$ alone, the
posterior has one peak (Proposition~\ref{prop:concave}), draws are exact, and the likelihood estimate is a lower bound. The two
forms are different laws with the same parameters, so one can start the fit of the other.

%======================================================================
\section{Evidence so far}
\label{sec:evidence}

\paragraph{Exactness.} In the test world (6 products, 3 categories, $r=2$) and all 11 model variants,
$\log p(S\mid z)$ from \eqref{eq:cond}--\eqref{eq:Z} equals enumeration of every basket to $10^{-9}$, and the
estimate \eqref{eq:is} with 4{,}000 draws equals enumeration plus $48\times48$ Gauss--Hermite quadrature of
\eqref{eq:marginal} to within 0.01 nats (\texttt{tests/test\_law.py}).

\paragraph{Synthetic data} (48 categories, 522 products). Held-out log-likelihood $-27.38$ against $-27.90$ for the
earlier form fully trained; mean generated size 7.92 against 8.69 in the data.

\paragraph{Real data} (64 held-out trips). The earlier form's old fit scores $-36.78$ with a peak-covering
estimate. The mixture form, in the full fit started from it:
\begin{center}\small
\begin{tabular}{lrrrr}\toprule
step & log-lik. & mean ESS & size, mean & size, 99\%\\\midrule
data & & & 7.26 & 42.8\\
240 & $-35.777$ & 0.35 & 7.66 & 41\\
360 & $\mathbf{-35.744}$ & 0.31 & 7.86 & 44.8\\
480 (whitening, $K=8$) & $-35.895$ & 0.27 & 8.23 & 46\\\bottomrule
\end{tabular}
\end{center}
The run was restarted from step 360 with the whitening estimate of Section~\ref{sec:whiten}; no generated basket
has reached the 120 cap at any point.

\section{Open points}
\begin{itemize}
\item The whitening step is approximate (Section~\ref{sec:whiten}).
\item The trained bound \eqref{eq:obj} with $K=8$ is below $\log p$; how much it biases the fitted parameters
  is measured only indirectly (validation at $K=256$).
\item In the real-data fit of Section~\ref{sec:evidence} only $\lambda,\phi,\psi,\rho_0,\rho_c,f$ were trained;
  price, promotion, competition and history terms kept the earlier fit's values. The fit from the additive
  base model, which trains them all, has not been run yet, and the price and promotion checks have to be
  re-run on this form.
\item The prior is one Gaussian. A mixture of a few Gaussians would keep Proposition~\ref{prop:concave} within each component.
\end{itemize}

\end{document}
